%%%%%%%%%%%%%%%%%%%%%%%%%%%%%%%%%%%%%%%%%%%%%%%%%%%%%%%%%%%%%%%%%%%%%%%%%%%%%%%%%%
%% This project aims to create a test template or exercise list from the        %%     
%% Federal University of Ceará (UFC).                                           %%
%% author: Maurício Moreira Neto - Doctoral student in Computer Science         %%
%% contacts:                                                                    %%
%%    e-mail: maumneto@ufc.br                                                   %%
%%    linktree:  https://linktr.ee/maumneto                                     %%
%%%%%%%%%%%%%%%%%%%%%%%%%%%%%%%%%%%%%%%%%%%%%%%%%%%%%%%%%%%%%%%%%%%%%%%%%%%%%%%%%%
\documentclass{ufcdocument}

\usepackage[utf8]{inputenc}
\usepackage[portuguese]{babel}
\usepackage{amssymb}
\usepackage{booktabs}
\usepackage[table]{xcolor}

\begin{document}

       %% Informations that will be insert in the table header 
\def\course{Lógica para Ciência da Computação}
\def\prof{Esdras Lins Bispo Jr.}
\def\semester{2026.1}
\def\codeCourse{ICE0239}
\def\registration{}
\def\student{}
\def\graduate{Bacharelado em Ciência da Computação}
\def\theme{Avaliação da Competência LCC C1.3:\\ \texttt{Construção de Tabelas-Verdade.}\\
07 de abril de 2026
}
    %% Table with the header
    \makeheader
    
    %% Space for the instructions
    % \fbox{
    %     \parbox{\textwidth}{
    %         \begin{minipage}{\textwidth}
    %             \makeinstructions
    %             {
    %                 \begin{instlist}
    %                     \item A avaliação é individual e não é pesquisada.
    %                     \item Preencha o cabeçalho da folha pergunta com seus dados.
    %                     \item  Todas as folhas respostas devem conter o nome a a matrícula do aluno.
    %                     \item O preenchimento das respostas deve ser feito utilizando caneta (preta ou azul).
    %                 \end{instlist}
    %             }
    %         \end{minipage}
    %     }
    % }
    %% Space between the instructions and the questions.
    \vspace{0.5cm}
    
     \begin{question}
        \item Construir a tabela-verdade da seguinte proposição
		\begin{center}
			$(\sim r \leftrightarrow p$ $\vee$ $q$) $\rightarrow$ $\sim q$
		\end{center}
    \end{question}
{
    \color{blue} {\bf Resposta:}
    \begin{center}
        \vspace{0.5cm}
        % Tabela-verdade para (\sim r \leftrightarrow p \vee q) \rightarrow \sim q
        \begin{tabular}{c|c|c||c|c|c|c||c|c}
            & & & & & (1) & (2) & (3) & \\
            \hline
            $p$ & $q$ & $r$ & $\sim r$ & $p \vee q$ & $\sim r \leftrightarrow (p \vee q)$ & $\sim q$ & (1) $\rightarrow$ (2) &
            \\
            \hline
            V & V & V & F & V & F & F & V \\
            V & V & F & V & V & V & F & F \\
            V & F & V & F & V & F & V & V \\
            V & F & F & V & V & V & V & V \\
            F & V & V & F & V & F & F & V \\
            F & V & F & V & V & V & F & F \\
            F & F & V & F & F & V & V & V \\
            F & F & F & V & F & F & V & V \\
            \hline
        \end{tabular}
    \end{center}
}
 
    \vfill

    \begin{center}
    {\small
    \setlength{\tabcolsep}{10pt}
    \renewcommand{\arraystretch}{2}
    \definecolor{LightGray}{gray}{0.92}
    \begin{tabular}{p{0.47\textwidth}cc}
    \toprule
    \addlinespace[2pt]
    \rowcolor{black}
    \multicolumn{3}{c}{\sffamily\fontsize{12pt}{14pt}\selectfont\textcolor{white}{Rubrica C1.3: Construção de Tabelas-Verdade.}} \\
    \addlinespace[2pt]
    \cmidrule(lr){1-3}
    \midrule
    \textbf{Critérios} & \textbf{Em formação} & \textbf{Proficiente} \\
    \midrule
    \rowcolor{LightGray} Apresentação da Tabela-Verdade & $\square$ & $\square$ \\
    Estruturação das Valorações & $\square$ & $\square$ \\
    \rowcolor{LightGray} Ordem de Precedência & $\square$ & $\square$ \\
    Semântica dos Conectivos & $\square$ & $\square$ \\
    \bottomrule
    \end{tabular}
    }
    \end{center}

\end{document}